% Shared module: notation table (one symbol, one meaning). Checked by tools/check_notation.py.
\begin{table}[htbp]
\centering\small
\caption{Notation for the analysis. Each symbol has one meaning throughout the paper; \cref{tab:notation-alg} lists the notation of the algorithm.}\label{tab:notation}
\begin{tabularx}{\linewidth}{@{}l X l@{}}
\toprule
Symbol & Meaning & Introduced\\
\midrule
\multicolumn{3}{@{}l}{\emph{Operators and their algebra}}\\
$\ring=\C[\tau,\xi]$ & ring of constant-coefficient operators; $\tau\leftrightarrow\partial_t$, $\xi\leftrightarrow\partial_x$ & \cref{def:ring}\\
$P,\ Q_j,\ P_i,\ e_i$ & an operator and its symbol; a factor; an irreducible factor and its multiplicity & \cref{sec:ring,sec:factor}\\
$\Var(P)$ & characteristic variety $\{P=0\}\subset\C^2$ & \cref{def:ring}\\
$\kappa,\ \lambda(\kappa)$ & real wavenumber ($\xi=i\kappa$); temporal rate of a mode & \cref{sec:ring}\\
$\Lambda_j(\xi)$ & branch of an evolution factor $Q_j=\tau-\Lambda_j(\xi)$ & \eqref{eq:lagrange}\\
$\Ann(u)$ & annihilator: all operators that kill $u$ & \cref{def:ann}\\
$\Gu,\ \Ju,\ \sigp(u)$ & intrinsic operator $\gcdop\Ann(u)$; cofactor ideal; point spectrum $\Var(\Ju)$ & \eqref{eq:GJ}\\
$\EP$ & exponential-polynomials & \cref{def:ep}\\
$N,\ \deltaP$ & number of factors; dimension of the overlap between their solution spaces & \cref{thm:factor}\\
$m_j,\ d_j$ & degree of factor $j$ in $\tau$ and in $\xi$ & \cref{sec:alg-loop}\\
\multicolumn{3}{@{}l}{\emph{Geometry and energy}}\\
$\sumop$ & sum operator $(u_1,\dots,u_N)\mapsto\sum_ju_j$ & \cref{sec:cond}\\
$\gN$ & conditioning constant of $\sumop$; error amplification $\gN^{-1/2}$ & \eqref{eq:gammaN}\\
$\cF,\ \thetaF$ & Friedrichs cosine and angle; $\gamma_2=1-\cF$ & \cref{sec:cond}\\
$\Hs$ & Sobolev space of index $s$ & \cref{sec:cond}\\
$K;\ \mathcal A_n,\ \flow_n,\ \Delta_n$ & number of additive terms of a factor; a term; its flow $e^{t\mathcal A_n}$; increment $\flow_n-\Id$ & \cref{sec:eint}\\
$\hat{\mathcal A}_n(\kappa)$ & Fourier symbol of the term $\mathcal A_n$ & \eqref{eq:eint-fourier}\\
$\Eint,\ \rho$ & interaction term; interaction ratio $\norm{\Eint}/\norm{u(T)-u_0}$ & \eqref{eq:eint-def}, \eqref{eq:energy}\\
$w_j$ & energy share of group $j$ & \eqref{eq:energy}\\
\bottomrule
\end{tabularx}
\end{table}

\begin{table}[htbp]
\centering\small
\caption{Notation for the algorithm (continued from \cref{tab:notation}).}\label{tab:notation-alg}
\begin{tabularx}{\linewidth}{@{}l X l@{}}
\toprule
Symbol & Meaning & Introduced\\
\midrule
$T,\ L$ & length of the time window; spatial period & \cref{sec:alg-interp}\\
$\Delta t,\ \Delta x;\ h$ & grid spacings; mean spacing $\sqrt{TL/n}$ of $n$ scattered points & \cref{sec:nyquist}\\
$\noise,\ \sigbin$ & observation noise standard deviation; its value in one Fourier bin & \cref{sec:alg-interp,sec:nyquist}\\
$\mathcal K,\ z$ & band of usable wavenumbers; detection threshold in units of $\sigbin$ & \cref{sec:nyquist}\\
$h_\kappa$ & spread $\max_{\ell,\ell'}\abs{\kappa_\ell-\kappa_{\ell'}}$ of the wavenumbers on which a factor is fitted & \cref{prop:counting}\\
$A_p,\ \sigma_{\lambda,p}$ & amplitude of mode $p$; standard deviation of its rate & \cref{sec:alg-pencil}\\
$\Nsamp,\ u^{(q)},\ \pi_q$ & number of posterior field samples; sample $q$; its weight & \cref{sec:alg-interp}\\
$\Mfit$ & a structure (number and degrees of the factors) & \cref{sec:alg-loop}\\
$B_j$ & orthonormal basis of the sampled kernel of $Q_j$ & \cref{sec:alg-interference}\\
$E_j$ & kernel exponentials of $Q_j$ evaluated at the observation points & \cref{sec:alg-accept}\\
$\omega^2$ & prior variance of the component coefficients & \eqref{eq:linmodel}\\
$\eta_r,\ \ell^{(r)}$ & amplitudes and length scales of the Gaussian-process kernel & \eqref{eq:gp}\\
$\eta,\ \eta_h,\ \pi_h,\ N_\eta$ & hyperparameters; a draw of them; its weight; number of draws & \eqref{eq:marginal-eta}\\
$\beta,\ \cloud$ & coefficient vector of a factor; cloud of factors & \eqref{eq:cloud}\\
$\alpha$ & concentration of the Dirichlet process & \cref{sec:bayes-cloud}\\
$k,\ \keff,\ \bar w_j$ & number of clusters (groups) found; effective number of groups; cluster masses & \cref{sec:bayes-cloud}\\
\FC, \IDK, \MIS, \UND & output labels & \cref{sec:alg-labels}\\
\bottomrule
\end{tabularx}
\end{table}
